\section{Analysis and Discussion}
\label{sec:discussion}

\paragraph{Why the methods perform as they do.} The results follow the two mechanisms of
Section~\ref{sec:theory} closely. For \emph{state tracking} what matters is which transitions the recurrence can
express, not how large the state is: all recurrent models had the same 8192-float state, and the outcome was
decided by the eigenvalue range and by whether the transition can be non-diagonal. For \emph{recall} what matters
is how much the state can hold and how efficiently it is used: attention never runs out of memory, the SSM breaks
at a point that moves with its state size, and the delta rule postpones the break by overwriting instead of
accumulating. No single recurrent model is best at both, and attention is best at recall but fails to
generalise on state tracking. This is the empirical case for the hybrid architectures used in practice.

\paragraph{Failure case 1: recall beyond state capacity.} The SSM with $N=16$ solves 16 pairs (0.991) but reaches
only 0.134 with 32 pairs, and the $N=64$ model fails at 64 pairs (0.053). The failure is abrupt rather than
gradual, and its position moves with the state size, which identifies capacity as the cause. A caveat: ``failed
within 20{,}000 steps'' bounds what was learned, not what is learnable, and one seed was used per configuration.
That the breaking point moves monotonically with $N$ supports the capacity explanation over an optimisation
accident.

\paragraph{Failure case 2: shortcuts that fit the training length.} Gated DeltaNet with $\beta\in(0,1)$ reaches
0.997 on parity at length 64 but only 0.586 at 512. By the proposition of Section~\ref{sec:math} its transitions
cannot implement a sign flip, so its near-perfect training-length accuracy must come from a solution that exploits
the limited length, for example by counting within a bounded window. The same holds for the Transformer on
$\mathbb{Z}_3$ (0.991 vs.\ 0.422). In-distribution accuracy alone would have suggested these models learned the
task; only the length-generalisation test reveals otherwise.

\paragraph{Failure case 3: expressible but not learned (rotational SSM on $\mathbb{Z}_3$).} The rotational SSM can
represent $\mathbb{Z}_3$ exactly; a unit test confirms that a rotation by $2\pi/3$ returns the state after three
steps. Yet training never finds this solution. It fits short sequences (1.000 at length 16), becomes unreliable at
length 64 (0.719, two of three seeds fail), and never generalises (at most 0.43 at 512). The learned weights
(Figure~\ref{fig:rotation-angles}) show why: the angles are not multiples of $120^\circ$, tokens 1 and 2 receive
nearly identical angles, and part of the memory decays, so the network built a short-range shortcut. We see two
reasons why gradient descent struggles here. First, the loss depends on the angle $\theta$ through terms like
$\cos(k\theta)$ for all $k\le T$, so the loss surface in $\theta$ oscillates more rapidly the longer the training
sequences are; this matches the drop in training-length accuracy from 1.000 to 0.719 as $T$ grows from 16 to 64.
Second, the target angle $2\pi/3$ is an interior value that must be hit exactly, because any error accumulates over
the sequence. In contrast, a sign flip is reached by pushing the decay towards $-1$, and a reflection by pushing
$\beta$ towards 2; these are saturating parametrisations that gradient descent approaches exponentially fast.

\paragraph{Failure case 4: $\mathbb{Z}_3$ resists every model, although $S_3$ does not.} DeltaProduct with
$\beta\in(0,2)$ generalises on $S_3$ (0.912 at length 512) but not on $\mathbb{Z}_3$ (0.511), even though
$\mathbb{Z}_3$ is a subgroup of $S_3$ (its rotations). We offer a tentative explanation in line with failure case 3.
In $S_3$, the three transpositions can be implemented by single exact reflections ($\beta\to 2$, saturating), and
the 3-cycles then arise automatically as products of two such reflections. In $\mathbb{Z}_3$, every non-identity
token must itself be a rotation by exactly $\pm120^\circ$, which requires the two keys of the token to be at
exactly $60^\circ$ to each other, an interior value again. Testing this hypothesis, for example by initialising
the keys at the required angle or by training with a length curriculum, is left for future work.

\paragraph{Methodology lesson: an insufficient training budget ranks models wrongly.} Our first recall sweep
(length 512, vocabulary 4096, 4000 steps) produced unusable results: even attention reached only 0.12--0.30
with 8 pairs, and accuracy was non-monotonic in the number of pairs (e.g.\ the $N=64$ SSM scored 0.41 with 8 pairs
but 0.99 with 16). Recall is learned in a sudden jump after a plateau, so a small fixed budget measures whether a
run happened to jump in time, not what the model can do. Checks on the CPU excluded two other explanations (a
short convolution for attention, a smaller vocabulary). A calibration run with up to 20{,}000 steps and early
stopping then gave every model at least 99\% on 16 pairs, with steps-to-90\% of 500 (attention + short conv), 1500
(Gated DeltaNet), 5500 (SSM) and 6500 (plain attention). The final recall study uses this budget.

\paragraph{Hyperparameter sensitivity.} Three sensitivities stood out. (i) A short causal convolution, which gives
each layer direct access to the previous few tokens, makes attention learn recall about 13$\times$ faster (500 vs.\
6500 steps); without it attention must first learn a separate previous-token mechanism. (ii) In a preliminary CPU
experiment, attention with one 64-dimensional head reached 91\% on a small recall task while the same model with
two or four narrower heads stalled near 53\%; we therefore used a single head for recall. (iii) The random seed
matters for the harder state-tracking cases: Gated DeltaNet with $\beta\in(0,2)$ on $S_3$ ranges from 0.28 to 0.83
at length 512 across seeds, and the rotational SSM solves parity on one seed out of three. Reporting three seeds
was therefore necessary; a single seed would have supported opposite conclusions.

\paragraph{Ablations.} The experiments double as ablations of single mechanisms with everything else fixed: the
eigenvalue range (SSM and DeltaNet, Table~\ref{tab:state}), rotational transitions, one versus two Householder steps
per token, the training length of the rotational SSM (Table~\ref{tab:rotation-length}), the state size of the SSM
(Table~\ref{tab:recall}) and the short convolution in attention.

\paragraph{Computational requirements.} Every model trains on a 6\,GB laptop GPU in minutes. At the lengths that
fit, however, the fused FlashAttention kernel \citep{dao2022flashattention} beats our linear-time layers, whose
asymptotic advantage would only show beyond roughly 32k tokens. The quality of the GPU kernel matters as much as
the complexity class, which is why production SSMs ship hand-written kernels. Our Triton kernel illustrates the
difficulty. It is faster and leaner than the PyTorch code but reaches only about 5\% of the GPU's memory
bandwidth. It is neither compute-bound nor bandwidth-bound but latency-bound: the loop over time is strictly
sequential and only $b\cdot H\cdot P/B_P = 64$ programs run in parallel, too few to hide memory latency on 20
streaming multiprocessors. This is precisely the motivation for Mamba-2's chunked algorithm, which turns most of
the scan into matrix multiplications that can use tensor cores; a Triton version of Algorithm~\ref{alg:scan} is
the natural next step. In generation, the recurrent models' constant state saves memory (32\,MB vs.\ 512\,KB per
layer at 8k tokens), but per-token latency at batch 4 is dominated by kernel-launch overhead.

\paragraph{Limitations.} The models are small (two layers, under half a million parameters) and the tasks are
synthetic, so the results characterise mechanisms rather than language-modelling quality. Recall results use one
seed. The Triton kernel implements only the forward pass, so training used the PyTorch chunked algorithm. The
benchmarks were taken on a laptop GPU whose clock varies with load and temperature; the scan table therefore
reports only lengths of at least 4k tokens, where timings were stable. Finally, our rotational SSM uses a
simplified version of the Mamba-3 parametrisation (angles from a separate projection, Euler-style input term), and
may behave differently from the original.
